\documentclass[12pt]{article}
\usepackage[bottom]{footmisc} 
\renewcommand{\thempfootnote}{\arabic{mpfootnote}}

\usepackage[hidelinks]{hyperref}

\usepackage[utf8]{inputenc}

\usepackage[T1]{fontenc}
\usepackage[ttdefault=true]{AnonymousPro}

\usepackage{multicol}
\setlength{\columnsep}{1cm}

\usepackage{abc}
\usepackage{graphicx}
\usepackage{musicography}

\usepackage[margin=0.5in]{geometry} 
\geometry{letterpaper}


\usepackage{fancyhdr} 
\pagestyle{fancy} % options: empty , plain , fancy
\renewcommand{\headrulewidth}{0pt} % customise the layout...
\lhead{}\chead{}\rhead{}
\lfoot{}\cfoot{\thepage}\rfoot{}


\begin{document}

\thispagestyle{empty}

\begin{abc}[name=stannes,options={-O= +c -F single.fmt}]
X:1
T:Saint Anne's Reel
M:C|
K:D
"D"fedf edcB | A2FA DAFA | "G"B2GB EBGB | "D"A2FA DAFA | 
fedf edcB | A2FA DAFA | "G"Beed "A"cABc |1 "D"d2d2d3e :|2 "D"d2d2d2ag | 
|:"D"fdfa fdfa | "G"aggf g2gf | "A"edcB Acea | "D"baa^g "A7"a2a=g | 
"D"fdfa fdfa | "G"aggf g2gf | "A"edcB Aceg |1 "D"f2d2d2 ag :|2 "D"f2d2d4 |
\end{abc}

\vskip 1.0in
\noindent\makebox[\textwidth][c]{%
  \begin{minipage}{7.0in}
    \setlength{\parindent}{1.5em}


The well-known tune, St. Anne's Reel, is believed to be French
Canadian, as the first recording came out in 1927 as \textit{Quadrille
  du loup garou - 4ème partie} by Montreal fiddler, Willie
Ringuette. The tune was recorded again in the next few years under the
titles, \textit{Reel des esquimaults} and \textit{Reel lune de miel}
by Joseph Allard and Isidore Soucy
respectively.\footnote{\url{https://tunearch.org/wiki/Annotation:St._Anne\%27s\_Reel}}
`St. Anne' most likely refers to the Sainte-Anne-de-Bellevue
municipality on the western end of Montreal.\footnote{Duval,
  Jean. \textbf{La Musique de Joseph Allard 1873-1947}, 2018, p. 72.}

St. Anne's Reel is very popular among American Old Time, Irish, and
Canadian players probably because of the tune's ability to be
elaborated to fit the needs of the situation (e.g., instrumental
improvisation or dance steps) while remaining identifiably St. Anne's
Reel. I believe the reason is twofold: melodic rhythm and harmonic
progression. St. Anne's Reel's melodic rhythmic motif of quarter- 6
eighth notes (measure 2 in this version) can (and will) be moved
around to other measures as the player decides. The IV-V-I harmonic
progression (the G-A-D chords in measure 7-8 and 15-17) can easily
elevate the simplest theme or reign in the most elaborate of
variations. Saint Anne's Reel is a tune that all levels of playing can
make their own.

    ---Kitty Steetle
    
  \end{minipage}
  \vspace*{\fill}
}

\end{document}
